\subsection{Geometry Parameterization and Reference Domain}
\label{subsec:ref_mapping_concept}
When the physical domain $\Omega(\bm{\mu})$ undergoes shape or boundary modifications as a function of geometric parameters $\bm{\mu}_{\text{geo}} \subset \bm{\mu}$, standard projection-based model reduction faces severe challenges because snapshot vectors corresponding to different parameters reside on different spatial meshes and varying coordinate support.

In our IGA framework, geometric changes are handled naturally by mapping the parameter-dependent physical domain $\Omega(\bm{\mu})$ from a fixed, parameter-independent reference parametric domain $\hat{\Omega} = (0, 1)^d$ via the NURBS geometric mapping:
\begin{equation}
    \bm{x}(\bm{\xi}; \bm{\mu}) = \bm{\Phi}(\bm{\xi}; \bm{\mu}) = \sum_{A=1}^{N_{\text{cp}}} R_A(\bm{\xi}) \, \bm{P}_A(\bm{\mu}), \quad \bm{\xi} \in \hat{\Omega},
    \label{eq:geo_mapping}
\end{equation}
where $\bm{P}_A(\bm{\mu})$ denotes the parameterized control point coordinates.

\subsection{Transformation of Bilinear Forms to the Reference Configuration}
\label{subsec:ref_transformation}
By pulling back the volume integrals from $\Omega(\bm{\mu})$ to the fixed reference domain $\hat{\Omega}$, the mass and stiffness bilinear forms become:
\begin{align}
    m(\bm{w}, \bm{u}; \bm{\mu}) &= \int_{\hat{\Omega}} \rho(\bm{\mu}) \, \bm{w}(\bm{\xi}) \cdot \bm{u}(\bm{\xi}) \, J(\bm{\xi}; \bm{\mu}) \, \mathrm{d}\hat{\Omega}, \\
    a(\bm{w}, \bm{u}; \bm{\mu}) &= \int_{\hat{\Omega}} \left( \nabla_{\bm{\xi}} \bm{w} : \bm{J}^{-\top} \right) : \mathbb{C}(\bm{\mu}) : \left( \nabla_{\bm{\xi}} \bm{u} : \bm{J}^{-\top} \right) J(\bm{\xi}; \bm{\mu}) \, \mathrm{d}\hat{\Omega},
\end{align}
where $\bm{J}(\bm{\xi}; \bm{\mu}) = \frac{\partial \bm{x}}{\partial \bm{\xi}}$ is the Jacobian matrix of the geometric map and $J(\bm{\xi}; \bm{\mu}) = \det(\bm{J}(\bm{\xi}; \bm{\mu}))$ is the Jacobian determinant.

Because the knot vectors and basis functions $R_A(\bm{\xi})$ on $\hat{\Omega}$ remain identical across varying geometric configurations, the control-point degrees of freedom $\hat{\bm{u}}(t; \bm{\mu})$ share a common algebraic index space. This eliminates the need for remeshing or spatial interpolation across parameter instances and creates a unified algebraic foundation for parameterized model reduction.
